% Part: first-order-logic
% Chapter: models-theories
% Section: expressing-relations

\documentclass[../../../include/open-logic-section]{subfiles}

\begin{document}

\olfileid[fr]{fol}{mat}{exr}

\olsection{Exprimer des relations dans une structure}

\begin{explain}
L'un des principaux usages des !!{formula}s consiste à exprimer des
propriétés et des relations dans !!a{structure}~$\Struct M$ à partir
des symboles primitifs du langage~$\Lang L$ de~$\Struct M$. Voici ce que cela
signifie. Le !!{domain} de~$\Struct M$ est un ensemble d'objets.
Les !!{constant}s, les !!{function}s et les !!{predicate}s sont
interprétés dans~$\Struct M$ respectivement par des objets de
$\Domain M$, des fonctions sur~$\Domain M$ et des relations sur
$\Domain M$. Par exemple, si $\Obj A^2_0$ appartient à~$\Lang L$,
alors $\Struct M$ lui attribue une relation
$R = \Assign{{\Obj A^2_0}}{M}$. La formule
$\Atom{\Obj A^2_0}{\Obj v_1, \Obj v_2}$ \emph{exprime} cette
relation en ce sens : si une assignation de variables~$s$ envoie
$\Obj v_1$ sur $a \in \Domain{M}$ et $\Obj v_2$ sur
$b \in \Domain M$, alors
\[
Rab \text{\quad si et seulement si\quad} \Sat{M}{\Atom{\Obj A^2_0}{\Obj v_1, \Obj v_2}}[s].
\]
Notons la nécessité de faire intervenir les assignations de
variables : nous ne pouvons pas simplement écrire
« $Rab$ si et seulement si $\Sat{M}{\Atom{\Obj A^2_0}{a, b}}$ »,
car $a$ et $b$ ne sont pas des symboles de notre langage, mais des
!!{element}s de~$\Domain{M}$.

Nous disposons non seulement des !!{formula}s atomiques, mais aussi
des connecteurs logiques et des quantificateurs pour les combiner.
Des !!{formula}s plus complexes peuvent ainsi définir des relations
qui ne figurent pas directement parmi les composantes de~$\Struct M$.
Nous voulons comprendre comment procéder et, plus précisément,
quelles relations sont définissables dans !!a{structure}.
\end{explain}

\begin{defn}
Soient $!A(\Obj v_1,\dots, \Obj v_n)$ !!a{formula} de~$\Lang L$
dont seules $\Obj v_1$,\dots, $\Obj v_n$ sont libres, et
$\Struct M$ !!a{structure} pour~$\Lang L$. La formule
$!A(\Obj v_1,\dots, \Obj v_n)$ \emph{exprime la relation}
$R \subseteq \Domain M^n$ si et seulement si
\[
Ra_1\dots a_n \text{\quad si et seulement si\quad} \Sat{M}{\Atom{!A}{\Obj
    v_1,\dots, \Obj v_n}}[s]
\]
pour toute assignation de variables~$s$ telle que
$s(\Obj v_i) = a_i$ ($i = 1, \dots, n$).
\end{defn}

\begin{ex}
Dans le modèle standard de l'arithmétique~$\Struct N$, la
!!{formula} $\Obj v_1 < \Obj v_2 \lor \eq[\Obj v_1][\Obj v_2]$
exprime la relation $\le$ sur~$\Nat$. La !!{formula}
$\eq[\Obj v_2][\Obj v_1']$ exprime la relation successeur,
c'est-à-dire la relation $R \subseteq \Nat^2$ telle que $Rnm$ si
$m$ est le successeur de~$n$. La formule
$\eq[\Obj v_1][\Obj v_2']$ exprime la relation prédécesseur.
Les !!{formula}s
$\lexists[\Obj v_3][(\eq/[\Obj v_3][\Obj 0] \land
  \eq[\Obj v_2][(\Obj v_1 + \Obj v_3)])]$ et
$\lexists[\Obj v_3][\eq[(\Obj v_1 + {\Obj v_3}')][\Obj v_2]]$
expriment toutes deux la relation $\Obj <$.
Le symbole de prédicat~$<$ est donc superflu dans le langage de
l'arithmétique : il peut y être défini.
\end{ex}

\begin{explain}
Cette idée n'est pas intéressante seulement pour des !!{structure}s
particulières. Elle l'est aussi, plus généralement, lorsque nous
utilisons un langage pour décrire un ou plusieurs modèles visés,
autrement dit lorsque nous considérons des théories. Celles-ci ne
prennent souvent comme symboles primitifs que quelques
!!{predicate}s, alors que bien d'autres relations jouent un rôle
important dans le domaine qu'elles décrivent. Si ces autres
relations peuvent être exprimées systématiquement à partir de celles
qui interprètent les !!{predicate}s primitifs du langage, nous
disons qu'elles sont \emph{définissables} dans ce langage.
\end{explain}

\begin{prob}
Trouver des !!{formula}s de~$\Lang L_A$ qui définissent les
relations suivantes :
\begin{enumerate}
\item $n$ est compris entre $i$ et $j$;
\item $n$ divise $m$ (autrement dit, $m$ est un multiple de~$n$);
\item $n$ est un nombre premier ($n>1$ et aucun nombre autre que $1$ et $n$
  ne divise~$n$).\footnote{Correction éditoriale : la source omet la condition qui exclut le nombre un.}
\end{enumerate}
\end{prob}

\begin{prob}
Supposons que la formule $!A(\Obj v_1, \Obj v_2)$ exprime la
relation $R \subseteq \Domain M^2$ dans !!a{structure}~$\Struct M$.
Trouver des formules qui expriment les relations suivantes :
\begin{enumerate}
\item la relation inverse $R^{-1}$ de~$R$;
\item le produit relatif $R \mid R$;
\end{enumerate}
Peut-on exprimer $R^+$, la clôture transitive de~$R$ ?
\end{prob}

\begin{prob}
Soit $\Lang{L}$ le langage qui ne contient que le symbole de
prédicat binaire $<$ (aucune autre !!{constant}, !!{function} ni
!!{predicate}, hormis bien sûr~$\eq$). Soit $\Struct{N}$ la
structure telle que $\Domain{N} = \Nat$ et
$\Assign{<}{N} = \Setabs{\tuple{n,m}}{n < m}$. Démontrer que :
\begin{enumerate}
\item $\{ 0 \}$ est définissable dans $\Struct{N}$;
\item $\{ 1 \}$ est définissable dans $\Struct{N}$;
\item $\{ 2 \}$ est définissable dans $\Struct{N}$;
\item pour chaque $n \in \Nat$, l'ensemble $\{ n \}$ est
  définissable dans $\Struct{N}$;
\item tout sous-ensemble fini de $\Domain{N}$ est définissable dans
  $\Struct{N}$;
\item tout sous-ensemble cofini de $\Domain{N}$ est définissable
  dans $\Struct{N}$ (où $X \subseteq \Nat$ est cofini si et
  seulement si $\Nat \setminus X$ est fini).
\end{enumerate}
\end{prob}


\end{document}
